\section{全书架构}

这份手稿围绕九项教学任务展开：
\begin{enumerate}[nosep]
  \item 第~\ref{ch:formal-language} 章与第~\ref{ch:dynamics} 章定义模型的语言。
  \item 第~\ref{ch:planning} 章解释为什么计划在结构上不可或缺。
  \item 第~\ref{ch:consciousness-biology} 章引入实证性的意识层，并把状态转变与生物学证据连接起来。
  \item 第~\ref{ch:emergence} 章解释持久框架如何从重复互动中形成，而不是来自某一次孤立选择。
  \item 第~\ref{ch:advanced-dynamics} 章把数学语言从轻量级的障碍模型升级为更丰富的状态空间视角。
  \item 第~\ref{ch:cases} 章用多个领域的案例把模型具体化。
  \item 第~\ref{ch:interventions} 章把分析转成设计模式。
  \item 第~\ref{ch:exercises} 章检验读者是否能够复用这一框架。
  \item 附录保留记号、模板，以及被清晰隔离出来的前沿假说。
\end{enumerate}